\documentclass[11pt]{article}
\usepackage[utf8]{inputenc}
\usepackage{amsmath,amssymb,amsthm}
\usepackage{enumerate}
\usepackage[margin=1in]{geometry}
\usepackage{hyperref}

\newtheorem{theorem}{Theorem}
\newtheorem{lemma}[theorem]{Lemma}
\newtheorem{corollary}[theorem]{Corollary}
\newtheorem{definition}[theorem]{Definition}
\newtheorem{remark}[theorem]{Remark}
\newtheorem{proposition}[theorem]{Proposition}

\newcommand{\lf}{\mathsf{lf}}
\newcommand{\nd}{\mathsf{nd}}
\newcommand{\cas}{\mathsf{cas}}
\newcommand{\rec}{\mathsf{rec}}
\newcommand{\niter}{\mathsf{niter}}
\newcommand{\eval}{\mathsf{eval}}
\newcommand{\reify}{\mathsf{reify}}
\newcommand{\codeTerm}{\mathsf{codeTerm}}
\newcommand{\decodeTerm}{\mathsf{decodeTerm}}
\newcommand{\codeReify}{\mathsf{codeReify}}
\newcommand{\diagCode}{\mathsf{diagCode}}
\newcommand{\diag}{\mathsf{diag}}
\newcommand{\thS}{\mathsf{thS}}
\newcommand{\thTerm}{\mathsf{thTerm}}
\newcommand{\eqTree}{\mathsf{eqTree}}
\newcommand{\transCode}{\mathsf{transCode}}
\newcommand{\swapCode}{\mathsf{swapCode}}
\newcommand{\codedSubst}{\mathsf{codedSubst}}
\newcommand{\substAt}{\mathsf{substAt}}
\newcommand{\Tree}{\mathsf{Tree}}
\newcommand{\Term}{\mathsf{Term}}
\newcommand{\Eq}{\mathsf{Eq}}
\newcommand{\dThS}{D_{\thS}}
\newcommand{\ValidProof}{\mathsf{ValidProof}}
\newcommand{\ThSResult}{\mathsf{ThSResult}}
\newcommand{\TrueEqCode}{\mathsf{TrueEqCode}}
\newcommand{\IsTermCode}{\mathsf{IsTermCode}}
\newcommand{\evalCode}{\mathsf{evalCode}}
\newcommand{\coreTree}{\mathsf{coreTree}}
\newcommand{\CoreInv}{\mathsf{CoreInv}}

\title{Formalizing G\"odel's Incompleteness in a Binary Tree Calculus:\\
  A Guard/Rose-Style Approach in Agda}
\author{}
\date{March 2026}

\begin{document}
\maketitle

\begin{abstract}
We describe an Agda formalization of G\"odel's incompleteness theorems
in a binary tree calculus inspired by Rose's system B* and Guard's 1963 lecture
notes. The formalization uses intrinsically scoped de Bruijn terms over binary
trees, with primitive recursion ($\rec$) and right-spine iteration ($\niter$) as
the computational primitives.

The main results include: (1)~a coded diagonal operator with a
correctness proof linking tree-level and syntactic substitution; (2)~an
internal theorem enumerator $\thS$ with reflexivity, symmetry, transitivity,
and computation axioms (cas-leaf, rec-leaf),
together with a full correctness proof; (3)~a semantic soundness theorem
via a validity predicate on proof trees ($\ValidProof$) and a decoded-equation
witness ($\ThSResult$); (4)~a proof-quoting operator $\dThS$ with metalevel
and internal correctness; (5)~a decode round-trip lemma
$\decodeTerm\;n\;(\codeTerm\;t) = \mathsf{just}\;t$; (6)~a G\"odel
sentence construction with both code-level and semantic fixed-point theorems;
(7)~G\"odel's First Incompleteness Theorem: a specific \emph{true} equation
($\niter\;\lf\;\lf\;(\mathsf{pair}\;v_1\;v_0) = \mathsf{leaf}$, both sides
evaluating to $\lf$) is proved to lie outside the range of $\thS$ via a novel
\emph{coreTree quotient} technique; and (8)~meta-level consistency and a
Rose-style Second Incompleteness Theorem: $\thS$ never produces the false
equation code, and the self-referential G\"odel consistency sentence is true
but unprovable.

All proofs use Agda's \texttt{--safe} flag with no postulates, no holes,
and no \texttt{with}-abstraction.
\end{abstract}

\section{The Object System}

\subsection{Binary Trees}

The sole data type is:
\[
  \Tree ::= \lf \mid \nd(a, b)
\]
This is Rose's system B*, where all computation operates on binary trees.

\subsection{Terms}

Terms are intrinsically scoped using de Bruijn indices:
\[
  \Term\;n ::= \mathsf{var}\;i \mid \mathsf{leaf} \mid \mathsf{pair}\;t\;u
    \mid \cas\;t\;u\;v \mid \rec\;t\;z\;s \mid \niter\;t\;\mathit{st}\;s
\]
where $n$ is the number of free variables, $\cas$ binds 2 variables in $v$,
$\rec$ binds 4 variables in $s$, and $\niter$ binds 2 variables in $s$.

The computation rules are:
\begin{align*}
  \cas\;\lf\;u\;v &\to u \\
  \cas\;(\nd\;a\;b)\;u\;v &\to v[b/0,\, a/1] \\
  \rec\;\lf\;z\;s &\to z \\
  \rec\;(\nd\;a\;b)\;z\;s &\to s[\rec\;b\;z\;s/0,\, \rec\;a\;z\;s/1,\, b/2,\, a/3] \\
  \niter\;\lf\;\mathit{st}\;s &\to \mathit{st} \\
  \niter\;(\nd\;a\;k)\;\mathit{st}\;s &\to \niter\;k\;(s[k/0,\,\mathit{st}/1])\;s
\end{align*}

\subsection{Evaluation and Coding}

Environment-based evaluation gives $\eval : \Term\;0 \to \Tree$.
Reification injects trees as closed terms: $\reify : \Tree \to \Term\;0$,
with the round-trip property $\eval(\reify\;t) = t$.

Coding maps terms to trees: $\codeTerm : \Term\;n \to \Tree$, using a
uniform $\nd\;\mathit{tag}\;\mathit{payload}$ encoding.

\section{Coded Substitution and the Diagonal}

\subsection{Coded Substitution}

The function $\codedSubst\;\mathit{repl}\;k\;\mathit{code}$ performs substitution
at the tree (code) level, replacing the $k$-th variable in the coded term
$\mathit{code}$ with $\mathit{repl}$.

\begin{theorem}[Coded Substitution Correctness]
For all $k$, $r : \Term\;0$, $t : \Term\;(n{+}1)$:
\[
  \codedSubst\;(\codeTerm\;r)\;k\;(\codeTerm\;t) = \codeTerm\;(\substAt\;k\;r\;t)
\]
\end{theorem}

\subsection{The Diagonal Operator}

\begin{definition}
$\codeReify : \Tree \to \Tree$ computes $\codeTerm(\reify\;t)$ directly on trees.
The diagonal operator is:
\[
  \diagCode\;c = \codedSubst_0\;(\codeReify\;c)\;c
\]
This is Guard's $\mathsf{sub}(i, i)$: substitute the numeral of $c$ into $c$ itself.
\end{definition}

\begin{theorem}[Diagonal Correctness]
\[
  \diagCode\;(\codeTerm\;t) = \codeTerm\;(\substAt_0\;(\reify\;(\codeTerm\;t))\;t)
\]
\end{theorem}

\subsection{The Fixed-Point Construction}

\begin{definition}
For any schema $A : \Term\;1$, the diagonalization is:
\[
  \diag\;A = \substAt_0\;(\reify\;(\codeTerm\;A))\;A
\]
\end{definition}

\begin{theorem}[Fixed-Point Identity]
\[
  \codeTerm\;(\diag\;A) = \diagCode\;(\codeTerm\;A)
\]
\end{theorem}

\section{The Decode Round-Trip}

Decoding inverts coding. The function $\decodeTerm\;n : \Tree \to
\mathsf{Maybe}\;(\Term\;n)$ recovers a term from its tree code.

\begin{theorem}[Decode Round-Trip]\label{thm:round-trip}
For all $t : \Term\;n$:
\[
  \decodeTerm\;n\;(\codeTerm\;t) = \mathsf{just}\;t
\]
\end{theorem}

\noindent The proof is by mutual structural induction on $t$,
using subsidiary lemmas for natural numbers
($\mathsf{decodeNat}\;(\mathsf{natCode}\;k) = \mathsf{just}\;k$) and
finite indices
($\mathsf{natToFin}\;n\;(\mathsf{finToNat}\;i) = \mathsf{just}\;i$).

\begin{corollary}[Evaluation of Coded Terms]\label{cor:evalCode}
Defining $\evalCode\;c = \mathsf{maybeMap}\;\eval\;(\decodeTerm\;0\;c)$, we have:
\[
  \evalCode\;(\codeTerm\;t) = \mathsf{just}\;(\eval\;t)
\]
\end{corollary}

\begin{corollary}[Code Injectivity on Evaluation]\label{cor:codeEvalEq}
If $\codeTerm\;t_1 = \codeTerm\;t_2$ (as trees), then $\eval\;t_1 = \eval\;t_2$.
\end{corollary}

\noindent This follows from the round-trip via:
$\mathsf{just}\;(\eval\;t_1) = \evalCode\;(\codeTerm\;t_1)
= \evalCode\;(\codeTerm\;t_2) = \mathsf{just}\;(\eval\;t_2)$,
then $\mathsf{just}$-injectivity.

\section{The Theorem Enumerator}

\subsection{Metalevel Definition}

The function $\thS : \Tree \to \Tree$ enumerates coded equations (trees of
the form $\nd\;L\;R$ representing equations $L = R$). Proof trees encode:

\begin{itemize}
  \item $\lf$: default equation ($\mathsf{leaf} = \mathsf{leaf}$)
  \item $\nd\;\lf\;\mathit{payload}$: reflexivity ($\codeReify(\mathit{payload}) = \codeReify(\mathit{payload})$)
  \item $\nd\;(\nd\;\lf\;\lf)\;\mathit{payload}$: symmetry of $\thS(\mathit{payload})$
  \item $\nd\;(\nd\;\lf\;(\nd\;\lf\;(\nd\;u\;v)))\;\mathit{payload}$:
    cas-leaf axiom $\cas\;\lf\;U\;V = U$ (where $u = \codeTerm\;U$, $v = \codeTerm\;V$)
  \item $\nd\;(\nd\;\lf\;(\nd\;(\nd\;\lf\;\lf)\;(\nd\;z\;s)))\;\mathit{payload}$:
    rec-leaf axiom $\rec\;\lf\;Z\;S = Z$ (where $z = \codeTerm\;Z$, $s = \codeTerm\;S$)
  \item $\nd\;(\nd\;(\nd\;a\;b)\;c)\;\mathit{payload}$: transitivity of
    $\thS(\nd\;(\nd\;a\;b)\;c)$ and $\thS(\mathit{payload})$,
    with a tree-equality check on the middle terms
\end{itemize}

\subsection{Internal Representation}

\begin{definition}
$\thTerm : \Term\;1$ is a closed term that computes $\thS$ internally.
It uses $\rec$ for structural recursion on the proof tree, nested $\cas$
for dispatching on the proof structure (including an \emph{axiom dispatch}
sub-term for the cas-leaf and rec-leaf branches), and an $\niter$-based
work-stack machine for the tree-equality check in the transitivity branch.
\end{definition}

\begin{theorem}[Internal Theorem Enumerator Correctness]
For all $t : \Tree$:
\[
  \eval_{\rho[t/0]}\;\thTerm = \thS\;t
\]
\end{theorem}

\section{Internal Tree Equality}

\subsection{The Challenge}

Tree equality requires simultaneous structural recursion on two arguments.
The $\rec$ primitive gives recursion on \emph{one} argument, and the results
for the children compare each child against the \emph{whole} second argument,
not against the corresponding child. A direct $\rec$-based approach fails.

\subsection{The Solution: Linearize + Niter Work-Stack}

We decompose tree equality into three phases:

\begin{enumerate}
  \item \textbf{Linearize}: $\rec$ on both trees to produce a right-leaning
    tree (the ``clock'') with right-spine length equal to the total node count.
  \item \textbf{Work-stack machine}: $\niter$ on the clock, maintaining a
    state $\nd\;\mathit{flag}\;\mathit{stack}$ where $\mathit{stack}$ is a list
    of pairs to compare. Each step pops a pair, checks leaf/node agreement,
    and pushes child pairs for node-node matches.
  \item \textbf{Extract flag}: a final $\cas$ to read the flag ($\lf$ for equal,
    $\nd\;\lf\;\lf$ for unequal).
\end{enumerate}

\begin{theorem}[Internal Tree Equality Correctness]
\[
  \mathsf{extractFlagMeta}\;(\mathsf{runNiter}\;(\mathsf{linearize}\;(\nd\;a\;b))\;(\nd\;\lf\;(\nd\;(\nd\;a\;b)\;\lf))) = \eqTree\;a\;b
\]
\end{theorem}

\section{Soundness of the Theorem Enumerator}

\subsection{Unconditional Soundness}

\begin{theorem}[nd-Headedness]\label{thm:nd-headed}
For all $y : \Tree$:
\[
  \eqTree\;(\thS\;y)\;\lf = \mathsf{falseT}
\]
\end{theorem}

\noindent This says every $\thS$ output is of the form $\nd\;L\;R$ (never $\lf$).
The proof handles symmetry via $\swapCode$ (which always returns $\nd$-form)
and transitivity via $\transCode$ (which returns $\nd$-form or
$\mathsf{defaultCode} = \nd\;(\nd\;\lf\;\lf)\;(\nd\;\lf\;\lf)$).

\subsection{The Validity Predicate}

The $\thS$ function is total on all trees, but only produces semantically
meaningful equation codes when the input is a well-formed proof tree.
In particular, the computation axiom branches (cas-leaf, rec-leaf) carry
term codes as data, which must actually be valid $\codeTerm$ images.

\begin{definition}[Valid Proof]
$\IsTermCode\;n\;c$ holds when $c = \codeTerm\;t$ for some $t : \Term\;n$.

$\ValidProof : \Tree \to \mathsf{Set}$ is an inductive predicate mirroring
the structure of $\thS$:
\begin{itemize}
  \item Base ($\lf$) and reflexivity ($\nd\;\lf\;\mathit{payload}$): always valid.
  \item Symmetry: valid if the sub-proof is valid.
  \item Cas-leaf: valid if $\IsTermCode\;0\;u$ and $\IsTermCode\;2\;v$.
  \item Rec-leaf: valid if $\IsTermCode\;0\;z$ and $\IsTermCode\;4\;s$.
  \item Transitivity: valid if both sub-proofs are valid.
  \item Default branches: always valid (produce $\mathsf{defaultCode}$).
\end{itemize}
\end{definition}

\subsection{The TrueEqCode Predicate}

\begin{definition}
\[
  \TrueEqCode\;c =
  \begin{cases}
    \mathsf{Unit} & \text{if } c = \lf \\
    \mathsf{trueEqMaybe}\;(\evalCode\;L)\;(\evalCode\;R) & \text{if } c = \nd\;L\;R
  \end{cases}
\]
where $\mathsf{trueEqMaybe}\;(\mathsf{just}\;a)\;(\mathsf{just}\;b) = \Eq\;a\;b$,
and all other cases give $\mathsf{Unit}$ (vacuous truth).
\end{definition}

\begin{remark}[Vacuous Truth and Transitivity]\label{rmk:vacuous}
$\TrueEqCode$ with vacuous truth is preserved by $\swapCode$ (symmetry)
but \emph{not} by $\transCode$ (transitivity). If one side of an equation
code fails to decode while the other succeeds, vacuous truth holds locally,
but composing two such codes via $\transCode$ can yield a goal
$\Eq\;a\;c$ that is unprovable. This is the fundamental reason why
soundness must be restricted to valid proof trees.
\end{remark}

\subsection{The ThSResult Witness}

\begin{definition}
$\ThSResult\;c$ is inhabited when $c = \nd\;(\codeTerm\;l)\;(\codeTerm\;r)$
for some $l, r : \Term\;0$ with $\eval\;l = \eval\;r$. This is strictly
stronger than $\TrueEqCode$ because it guarantees both sides decode
(no vacuous truth).
\end{definition}

\begin{lemma}[Swap Preserves ThSResult]
If $\ThSResult\;c$ then $\ThSResult\;(\swapCode\;c)$.
\end{lemma}

\begin{lemma}[TransCode Preserves ThSResult]\label{lem:trans-thsr}
If $\ThSResult\;e_1$ and $\ThSResult\;e_2$ then $\ThSResult\;(\transCode\;e_1\;e_2)$.
\end{lemma}

\noindent The transitivity case works because $\ThSResult$ guarantees
both sides are $\codeTerm$ images. When $\eqTree$ matches the middle
terms (right of $e_1$ equals left of $e_2$ as trees), Corollary~\ref{cor:codeEvalEq}
gives the evaluation bridge: $\eval\;r_1 = \eval\;l_2$, which composes
with the sub-proofs via $\mathsf{eqTrans}$.

\subsection{The Main Soundness Theorem}

\begin{theorem}[Soundness of thS]\label{thm:soundness}
For all $p : \Tree$:
\[
  \ValidProof\;p \;\to\; \ThSResult\;(\thS\;p)
\]
and consequently:
\[
  \ValidProof\;p \;\to\; \TrueEqCode\;(\thS\;p)
\]
\end{theorem}

\noindent The proof is by induction on $\ValidProof\;p$:
\begin{itemize}
  \item \emph{Reflexivity}: $\thS$ produces $\nd\;(\codeReify\;a)\;(\codeReify\;a)$.
    By $\codeReify\text{-correct}$, this equals
    $\nd\;(\codeTerm\;(\reify\;a))\;(\codeTerm\;(\reify\;a))$, and $\eval$ of
    both sides is the same (reflexivity).
  \item \emph{Cas-leaf}: with $\IsTermCode$ witnesses $u = \codeTerm\;U$ and
    $v = \codeTerm\;V$, the output is
    $\nd\;(\codeTerm\;(\cas\;\mathsf{leaf}\;U\;V))\;(\codeTerm\;U)$.
    By the computation rule $\cas\;\lf\;U\;V \to U$, both sides evaluate
    to $\eval\;U$, so the proof is $\mathsf{refl}$.
  \item \emph{Rec-leaf}: analogous, using $\rec\;\lf\;Z\;S \to Z$.
  \item \emph{Symmetry}: by the inductive hypothesis and Swap preservation.
  \item \emph{Transitivity}: by the inductive hypothesis and Lemma~\ref{lem:trans-thsr}.
  \item \emph{Default branches}: produce $\mathsf{defaultCode} = \nd\;(\codeTerm\;\mathsf{leaf})\;(\codeTerm\;\mathsf{leaf})$,
    trivially true.
\end{itemize}

\section{The Proof-Quoting Operator $\dThS$}

\begin{definition}
The metalevel proof-quoting operator is:
\[
  \dThS\;p = \nd\;\lf\;(\thS\;p)
\]
The internal term is $\mathsf{dThSTerm} = \mathsf{pair}\;\mathsf{leaf}\;\thTerm : \Term\;1$.
\end{definition}

\noindent This uses the reflexivity branch of $\thS$: the proof tree
$\nd\;\lf\;\mathit{payload}$ proves $\codeReify(\mathit{payload}) = \codeReify(\mathit{payload})$.
Setting $\mathit{payload} = \thS\;p$, we obtain a proof that $\thS(p) = \thS(p)$.

\begin{theorem}[D\_thS Correctness --- Metalevel]
\[
  \thS\;(\dThS\;p) = \nd\;(\codeReify\;(\thS\;p))\;(\codeReify\;(\thS\;p))
\]
\end{theorem}
\begin{proof}
Definitional (by the reflexivity case of $\thS$).
\end{proof}

\begin{theorem}[D\_thS Correctness --- Internal]
\[
  \eval_{\rho[p/0]}\;\mathsf{dThSTerm} = \dThS\;p
\]
\end{theorem}
\begin{proof}
By $\mathsf{eqCong}\;(\nd\;\lf)\;(\mathsf{thTerm\text{-}correct}\;p)$.
\end{proof}

\begin{proposition}[D\_thS Validity and Soundness]
\mbox{}
\begin{enumerate}[(i)]
  \item $\ValidProof\;(\dThS\;p)$ for all $p$.
  \item $\TrueEqCode\;(\thS\;(\dThS\;p))$ for all $p$.
\end{enumerate}
\end{proposition}
\begin{proof}
(i) By $\mathsf{vpRefl}\;(\thS\;p)$.
(ii) By $\mathsf{trueEq\text{-}refl}\;(\thS\;p)$.
\end{proof}

\section{The G\"odel Sentence}

\begin{definition}
For a target tree $\mathit{target}$:
\begin{align*}
  \mathsf{godelSchema}\;\mathit{target} &= \mathsf{matchSub}\;\mathit{target}\;\thTerm \\
  \mathsf{godelSentence}\;\mathit{target} &= \diag\;(\mathsf{godelSchema}\;\mathit{target})
\end{align*}
where $\mathsf{matchSub}\;\mathit{target}\;t$ checks whether $\eval\;t$ equals
$\mathit{target}$ by structural decomposition.
\end{definition}

\begin{theorem}[G\"odel Code-Level Fixed Point]
\[
  \codeTerm\;(\mathsf{godelSentence}\;\mathit{target})
  = \diagCode\;(\codeTerm\;(\mathsf{godelSchema}\;\mathit{target}))
\]
\end{theorem}

\begin{theorem}[G\"odel Semantic Fixed Point]
\[
  \eval\;(\mathsf{godelSentence}\;\mathit{target})
  = \eqTree\;(\thS\;(\codeTerm\;(\mathsf{godelSchema}\;\mathit{target})))\;\mathit{target}
\]
\end{theorem}

\begin{theorem}[G\"odel Sentence Evaluation]
For $\mathit{target} = \lf$:
\[
  \eval\;(\mathsf{godelSentence}\;\lf) = \mathsf{falseT}
\]
\end{theorem}
\begin{proof}
By the semantic fixed point and Theorem~\ref{thm:nd-headed}: $\thS$ always
returns an $\nd$-headed tree, so $\eqTree\;(\thS\;c)\;\lf = \mathsf{falseT}$.
\end{proof}

\section{The CoreTree Quotient and G\"odel I}

\subsection{The CoreTree Function}

The key technical innovation is a tree-level function $\coreTree : \Tree \to \Tree$
that strips $\cas$-leaf and $\rec$-leaf code wrappers, acting as a normal form
for the equational theory generated by the computation axioms.

\begin{definition}
$\coreTree$ is defined by mutual recursion:
\begin{itemize}
  \item $\coreTree\;\lf = \lf$
  \item $\coreTree\;(\nd\;\mathit{tag}\;\mathit{payload}) = \mathsf{coreDispatch}\;\mathit{tag}\;\mathit{payload}$
\end{itemize}
where $\mathsf{coreDispatch}$ dispatches on the tag (with exhaustive,
non-overlapping pattern matching for \texttt{--exact-split}):
\begin{itemize}
  \item If $\mathit{tag} = \mathsf{tagCase}$: inspect the payload.
    If the payload is $\nd\;(\nd\;\lf\;\lf)\;(\nd\;c_U\;c_V)$ (i.e., the
    scrutinee code is $\codeTerm\;\mathsf{leaf} = \nd\;\lf\;\lf$), return
    $\coreTree\;c_U$. Otherwise return the code unchanged.
  \item If $\mathit{tag} = \mathsf{tagRec}$: analogously, if the scrutinee
    code is $\nd\;\lf\;\lf$, return $\coreTree\;c_Z$. Otherwise unchanged.
  \item All other tags: return $\nd\;\mathit{tag}\;\mathit{payload}$ unchanged.
\end{itemize}
\end{definition}

\noindent The crucial property: $\coreTree$ computes the same value on both
sides of any equation that $\thS$ can produce.

\subsection{The CoreInv Predicate}

\begin{definition}
\[
  \CoreInv\;c =
  \begin{cases}
    \mathsf{Unit} & \text{if } c = \lf \\
    \Eq\;(\coreTree\;L)\;(\coreTree\;R) & \text{if } c = \nd\;L\;R
  \end{cases}
\]
\end{definition}

\begin{lemma}[CoreInv Preserved by Swap]
If $\CoreInv\;c$ then $\CoreInv\;(\swapCode\;c)$.
\end{lemma}
\begin{proof}
$\swapCode\;(\nd\;L\;R) = \nd\;R\;L$. Apply $\mathsf{eqSym}$.
\end{proof}

\begin{lemma}[CoreInv Preserved by TransCode]\label{lem:coreinv-trans}
If $\CoreInv\;e_1$ and $\CoreInv\;e_2$ then $\CoreInv\;(\transCode\;e_1\;e_2)$.
\end{lemma}
\begin{proof}
When $e_1 = \nd\;L_1\;R_1$, $e_2 = \nd\;L_2\;R_2$, and the middle matches
($\eqTree\;R_1\;L_2 = \lf$, so $R_1 = L_2$ by $\eqTree$-sound), the result
is $\nd\;L_1\;R_2$. The proof chains:
\[
  \coreTree\;L_1 = \coreTree\;R_1 = \coreTree\;L_2 = \coreTree\;R_2
\]
where the middle step uses $\mathsf{eqCong}\;\coreTree$ applied to $R_1 = L_2$.
When the middle does not match, $\transCode$ returns $\mathsf{defaultCode}$,
which trivially satisfies $\CoreInv$ (both sides are $\nd\;\lf\;\lf$).
\end{proof}

\begin{lemma}[CoreInv for Computation Axioms]\mbox{}
\begin{enumerate}[(i)]
  \item $\CoreInv\;(\mathsf{casLeafCode}\;u\;v) = \Eq\;(\coreTree\;u)\;(\coreTree\;u) = \mathsf{refl}$.
  \item $\CoreInv\;(\mathsf{recLeafCode}\;z\;s) = \Eq\;(\coreTree\;z)\;(\coreTree\;z) = \mathsf{refl}$.
\end{enumerate}
\end{lemma}
\begin{proof}
For (i): $\mathsf{casLeafCode}\;u\;v = \nd\;(\nd\;\mathsf{tagCase}\;(\nd\;(\nd\;\lf\;\lf)\;(\nd\;u\;v)))\;u$.
The left side's code has tag $\mathsf{tagCase}$ with scrutinee $\nd\;\lf\;\lf = \codeTerm\;\mathsf{leaf}$.
By definition, $\coreTree$ strips this to $\coreTree\;u$, which equals $\coreTree\;u$ (the right side).
This is definitional in Agda (proved by \texttt{refl}). Part~(ii) is analogous.
\end{proof}

\subsection{The Main Invariant}

\begin{theorem}[CoreInv for All Trees]\label{thm:coreinv-all}
For all $y : \Tree$:
\[
  \CoreInv\;(\thS\;y)
\]
\end{theorem}
\begin{proof}
By structural induction on $y$, following the case structure of $\thS$:
\begin{itemize}
  \item Base and default branches: $\thS$ returns $\mathsf{defaultCode}$ or a
    reflexivity code. Both sides equal, so $\CoreInv$ holds by $\mathsf{refl}$.
  \item Reflexivity: $\thS$ returns $\nd\;(\codeReify\;p)\;(\codeReify\;p)$.
    $\mathsf{refl}$.
  \item Symmetry: by the inductive hypothesis and swap preservation.
  \item Cas-leaf / rec-leaf: by computation (definitional $\mathsf{refl}$).
  \item Transitivity: by Lemma~\ref{lem:coreinv-trans} with both inductive hypotheses.
\end{itemize}
Note: this theorem requires \emph{no} $\ValidProof$ hypothesis. The invariant
holds for \emph{all} trees because every branch of $\thS$ preserves it structurally.
\end{proof}

\subsection{Generalized Range Exclusion}

\begin{theorem}[Core-Unprovable]\label{thm:core-unprov}
For all $L, R : \Tree$, if $\coreTree\;L \neq \coreTree\;R$, then for all $y : \Tree$:
\[
  \thS\;y \neq \nd\;L\;R
\]
\end{theorem}
\begin{proof}
If $\thS\;y = \nd\;L\;R$, then by Theorem~\ref{thm:coreinv-all},
$\coreTree\;L = \coreTree\;R$, contradicting the hypothesis.
\end{proof}

\noindent This is a powerful metatheorem: it gives an effective, decidable
criterion for unprovability. Given any equation code $\nd\;L\;R$, compute
$\coreTree\;L$ and $\coreTree\;R$; if they differ, the equation is
\emph{never} produced by $\thS$, regardless of input.

\subsection{G\"odel I: True-but-Unprovable}

\begin{definition}
The \emph{niter-leaf equation}: $\niter\;\lf\;\lf\;(\mathsf{pair}\;v_1\;v_0) = \mathsf{leaf}$.
\end{definition}

\begin{theorem}[Niter-Leaf Equation is True]
$\eval\;(\niter\;\lf\;\lf\;(\mathsf{pair}\;v_1\;v_0)) = \eval\;\mathsf{leaf} = \lf$.
\end{theorem}
\begin{proof}
$\niter$ with clock $\lf$ returns the state immediately: $\niter\;\lf\;\mathit{st}\;s \to \mathit{st}$.
Since $\mathit{st} = \lf$, both sides evaluate to $\lf$. (Proved by $\mathsf{refl}$.)
\end{proof}

\begin{theorem}[Niter-Leaf Equation is Unprovable]\label{thm:godel1}
For all $y : \Tree$:
\[
  \thS\;y \neq \nd\;(\codeTerm\;(\niter\;\lf\;\lf\;(\mathsf{pair}\;v_1\;v_0)))\;(\codeTerm\;\mathsf{leaf})
\]
\end{theorem}
\begin{proof}
By Theorem~\ref{thm:core-unprov}. The left side's code has tag $\mathsf{tagNiter}$;
the right side's code has tag $\mathsf{tagLeaf} = \lf$. Neither is $\mathsf{tagCase}$
or $\mathsf{tagRec}$, so $\coreTree$ returns both unchanged. Tag discrimination:
$\mathsf{tagNiter} = \nd\;\lf\;(\nd\;\lf\;(\nd\;\lf\;(\nd\;\lf\;(\nd\;\lf\;\lf))))
\neq \lf = \mathsf{tagLeaf}$.
\end{proof}

\noindent This is G\"odel's First Incompleteness Theorem in a strong form: a specific
true equation lies outside the range of $\thS$ for \emph{all} inputs, not just
valid proof trees.

\section{Meta-Level Consistency}

\begin{theorem}[Meta-Level Consistency]\label{thm:consistency}
For all $y : \Tree$:
\[
  \thS\;y \neq \nd\;(\codeTerm\;\mathsf{leaf})\;(\codeTerm\;(\mathsf{pair}\;\mathsf{leaf}\;\mathsf{leaf}))
\]
\end{theorem}
\begin{proof}
By Theorem~\ref{thm:core-unprov}. The left code $\codeTerm\;\mathsf{leaf} =
\nd\;\lf\;\lf$ has $\coreTree$ value $\nd\;\lf\;\lf$ (tag $\lf = \mathsf{tagLeaf}$).
The right code $\codeTerm\;(\mathsf{pair}\;\mathsf{leaf}\;\mathsf{leaf})$ has
tag $\mathsf{tagPair} = \nd\;\lf\;(\nd\;\lf\;\lf) \neq \lf$.
\end{proof}

\noindent This establishes the consistency of $\thS$ as a metatheorem:
the false equation ``$\mathsf{leaf} = \mathsf{pair}\;\mathsf{leaf}\;\mathsf{leaf}$''
is never produced by $\thS$, regardless of input. This is strictly stronger than
the soundness-based consistency proof (which requires $\ValidProof$) because
it holds for \emph{all} trees.

\section{G\"odel II: The Consistency Sentence is Unprovable}

\subsection{The Consistency Sentence}

Recall that $\mathsf{godelSentence}\;\lf$ evaluates to $\mathsf{falseT}$
(Theorem~8.3), and the equation
\[
  \mathsf{godelSentence}\;\lf = \mathsf{reifyFalse}
\]
is true (both sides evaluate to $\mathsf{falseT}$). This equation is a
consistency-type statement: it says that $\thS$, applied to a specific
self-referential input (the code of the G\"odel schema), does not produce $\lf$
(the degenerate value).

\subsection{Unprovability of the Consistency Sentence}

\begin{theorem}[G\"odel II, Rose-Style]\label{thm:godel2}
For all $y : \Tree$:
\[
  \thS\;y \neq \nd\;(\codeTerm\;(\mathsf{godelSentence}\;\lf))\;(\codeTerm\;\mathsf{reifyFalse})
\]
\end{theorem}
\begin{proof}
By Theorem~\ref{thm:core-unprov}. The key structural observation:
\begin{itemize}
  \item $\mathsf{godelSentence}\;\lf$ is a $\cas$-headed term (because
    $\mathsf{matchSub}$ always wraps its scrutinee in $\cas$, and
    $\diag$ preserves the outermost constructor). Therefore
    $\codeTerm\;(\mathsf{godelSentence}\;\lf)$ has tag $\mathsf{tagCase}$.
  \item The scrutinee of this $\cas$ (after diagonalization) is
    $\rec\;(\reify\;c)\;\ldots$, a $\rec$-headed term. Its code has tag
    $\mathsf{tagRec} \neq \lf$, so $\coreTree$ does \emph{not} strip the
    cas-leaf wrapper. The tag remains $\mathsf{tagCase}$.
  \item $\mathsf{reifyFalse} = \mathsf{pair}\;\mathsf{leaf}\;\mathsf{leaf}$
    has code with tag $\mathsf{tagPair}$. Since this is neither $\mathsf{tagCase}$
    nor $\mathsf{tagRec}$, $\coreTree$ returns it unchanged.
  \item Tag discrimination: $\mathsf{tagCase} = \nd\;\lf\;(\nd\;\lf\;(\nd\;\lf\;\lf))
    \neq \nd\;\lf\;(\nd\;\lf\;\lf) = \mathsf{tagPair}$.
\end{itemize}
\end{proof}

\begin{remark}[The CoreTree Technique vs.\ Classical Approaches]
The classical proof of G\"odel II via the L\"ob/Hilbert--Bernays argument
requires derivability conditions D1--D3 and an internal provability predicate.
Our proof is purely \emph{structural}: the $\coreTree$ quotient gives an effective
range characterization of $\thS$ that makes both incompleteness theorems
immediate consequences of tag discrimination. The tradeoff is that our G\"odel~II
proves unprovability of a specific self-referential consistency sentence, rather
than the abstract ``if Con is provable then $\bot$'' form.
\end{remark}

\section{What Remains}

\subsection{Stronger G\"odel II}

The current G\"odel II uses $\mathsf{godelSentence}\;\lf$, whose consistency
content is that $\thS$ at one self-referential input does not produce $\lf$.
A stronger version would use $\mathit{target} = \mathsf{falseEqCode}$
(the code of $\mathsf{leaf} = \mathsf{pair}\;\mathsf{leaf}\;\mathsf{leaf}$),
making the consistency content explicit: ``$\thS$ at a self-referential input
does not prove the false equation.'' The proof is structurally identical
(same tag discrimination) but involves larger Agda normalizations.

\subsection{Classical G\"odel II via D1--D3}

For the full L\"ob-style argument, one would need:
\begin{itemize}
  \item An internal provability predicate (existential quantification over
    proof trees), requiring a formula layer beyond the equational setting.
  \item D2 (distributing provability over implication) and D3 (provability
    of provability), which in the equational setting correspond to internal
    composition of proof trees via $\transCode$.
  \item The L\"ob fixed-point argument composing these conditions.
\end{itemize}
A separate formalization (GuardBinaryGodel2.agda, 234 lines) carries this out
in a predicate-logic setting over the same binary tree arithmetic.

\subsection{Additional Computation Axioms}

The current $\thS$ has cas-leaf and rec-leaf. For a more complete system:
\begin{itemize}
  \item \textbf{cas-pair}: $\cas\;(\mathsf{pair}\;(\reify\;a)\;(\reify\;b))\;\mathsf{leaf}\;v
    = \reify\;(\eval_{\rho[a,b]}\;v)$.
  \item \textbf{rec-pair}: analogous, with 4-variable substitution on the RHS.
  \item \textbf{niter-leaf / niter-pair}: iteration base and step.
\end{itemize}
Adding these would require extending $\coreTree$ to handle the new axiom patterns,
but the quotient technique generalizes: each new axiom $f(\ldots) = g(\ldots)$
contributes a new stripping rule to $\coreTree$, and the $\CoreInv$ preservation
proofs follow the same structure.

\section{Design Principles}

\begin{itemize}
  \item \textbf{ASCII only}: No Unicode in Agda files.
  \item \textbf{Semi-explicit style}: Implicit arguments for types and indices,
    explicit for computational content.
  \item \textbf{No with-abstraction}: All case analysis via $\mathsf{eqSubst}$,
    $\mathsf{eqCong}$, $\mathsf{eqTrans}$, or helper functions.
  \item \textbf{--safe --without-K --exact-split}: Maximum safety guarantees.
  \item \textbf{Guard-style, not Rose-style}: The fixed-point construction
    follows Guard's ``diagonalize the code'' approach rather than Rose's
    universal internal interpreter approach, avoiding the index-shifting
    obstruction that blocked earlier attempts at internalizing $\codedSubst$.
  \item \textbf{Validity-restricted soundness}: Semantic soundness
    ($\TrueEqCode$) is proved only for valid proof trees ($\ValidProof$),
    not for all trees. This resolves the vacuous-truth compositionality
    obstruction (Remark~\ref{rmk:vacuous}).
\end{itemize}

\section{Statistics}

The Rose formalization comprises 13 Agda files:
\begin{itemize}
  \item 10 core infrastructure files (Base through Internal)
  \item 3 higher-level files (DiagCode, FixedPoint, ThInt, Godel)
  \item 0 postulates, 0 holes, \texttt{--safe --without-K --exact-split}
  \item 36+ unit tests (all verified by \texttt{refl})
\end{itemize}

\noindent Key theorems proved:
\begin{enumerate}
  \item Coded substitution correctness ($\codedSubst$-correct)
  \item Diagonal correctness ($\diagCode$-correct)
  \item Fixed-point identity ($\diag$-code)
  \item Internal tree equality correctness ($\eqTree$-check-main)
  \item Internal theorem enumerator correctness ($\thTerm$-correct)
  \item Decode round-trip ($\decodeTerm\;n\;(\codeTerm\;t) = \mathsf{just}\;t$)
  \item Full soundness ($\ValidProof \to \ThSResult$)
  \item $\dThS$ correctness (metalevel and internal)
  \item G\"odel sentence evaluation and fixed-point
  \item G\"odel I: true-but-unprovable equation (Theorem~\ref{thm:godel1})
  \item Meta-level consistency (Theorem~\ref{thm:consistency})
  \item G\"odel II: consistency sentence unprovable (Theorem~\ref{thm:godel2})
\end{enumerate}

\end{document}
